\makeatletter
\def\input@path{{./}{../}{../../}{../../../}{../../../../}{preamble/}{../preamble/}{../../preamble/}{../../../preamble/}}
\makeatother

\input{preamble.tex}
\input{graphs/highlights.tex}
% \GTenHighlightsfalse % Uncomment to disable instructional highlights.
\input{graphs/specific_highlights.tex}
\COneTonePracticePalette


\input{graphs/G11_C6C7_graphs.tex}
\input{graphs/G11_Histograms.tex}
\usepackage{longtable}

\setlength{\DistributionGraphWidth}{0.45\textwidth}
\setlength{\DistributionGraphHeight}{4cm}

\title{Grade 11 -- Term 1 \\ C6--C7 Practice Activity with Solutions}
\author{}
\date{}

\pagestyle{empty}

% ============================================================
% GRAPH WRAPPER
% ============================================================

\begin{document}

	\maketitle
	\setcounter{tocdepth}{1}
	\setcounter{secnumdepth}{2}
	\phantomsection
	\label{toc}
	\tableofcontents
	
	\section*{Evaluated criterion}
	\begin{itemize}
		\item [\textbf{C6.}] Estimates interval probabilities using relative frequencies from generated samples and compares them with theoretical probabilities.
		
		\item [\textbf{C7.}] Calculates sample means and variances for different sample sizes and compares them with the theoretical mean and variance.
		
		
	\end{itemize}
	
	\vspace{-4mm}
	
% ==========================================================================
% C6 INTRODUCTION
% ==========================================================================

\section{Introduction}
\label{sec:c6-introduction}
\Needspace{8\baselineskip}
\subsection*{Mathematical highlight guide}
\phantomsection\label{tone:C6C7:legend}
\PracticeHighlightLegend{G11-T1-C6-C7}


\SectionProblemsTableOfContents{
	\SectionProblemLink{sec:c6-intro-samples-histograms}{Generated samples, sample sizes, and histograms} \\
	\SectionProblemLink{sec:c6-intro-distributions}{Continuous distributions, theoretical and empirical probabilities} \\
	\SectionProblemLink{sec:c6-intro-sample-mean-variance}{Sample mean and variance from generated data}
}
\vspace{8mm}

Compare generated samples with a theoretical distribution using interval
probabilities (C6), sample means, and sample variances (C7).
Theoretical probabilities are areas under the \CtwoToneDensity{probability} \CtwoToneDensity{density} \CtwoToneDensity{function}
(PDF); \CsixToneEmpirical{empirical} \CsixToneEmpirical{probabilities} are \CsixToneEmpirical{relative} \CsixToneEmpirical{frequencies} in the sample.


% ==========================================================================
% CONTINUOUS PROBABILITY DISTRIBUTIONS
% ==========================================================================
\subsection{Real-world data, probability distributions, and histograms}
\label{sec:c6-intro-samples-histograms}

A probability distribution models variation in a variable. A sample contains
$\CsixToneCount{n}$ collected or generated observations:

$$
\ConeToneVariable{x_1},\ConeToneVariable{x_2},\ldots,\ConeToneVariable{x_n}.
$$

A histogram groups observations into intervals and shows their frequencies
for comparison with the theoretical distribution.

\begin{center}
	\begin{tabular}{ccc}
		\begin{minipage}[t]{0.30\textwidth}
			\centering
			\CSampleCountHistogramTwenty
			\captionof{figure}{
				Count histogram from a sample of $\CsixToneCount{n}=\CsixToneCount{20}$ observations.
			}
			\label{fig:c6-sample-count-histogram-20}
		\end{minipage}
		&
		\begin{minipage}[t]{0.30\textwidth}
			\centering
			\CSampleCountHistogramHundred
			\captionof{figure}{
				Count histogram from a sample of $\CsixToneCount{n}=\CsixToneCount{100}$ observations.
			}
			\label{fig:c6-sample-count-histogram-100}
		\end{minipage}
		&
		\begin{minipage}[t]{0.30\textwidth}
			\centering
			\CSampleCountHistogramThousand
			\captionof{figure}{
				Count histogram from a sample of $\CsixToneCount{n}=\CsixToneCount{1000}$ observations.
			}
			\label{fig:c6-sample-count-histogram-1000}
		\end{minipage}
	\end{tabular}
\end{center}

\begin{center}
	\begin{tabular}{ccc}
		\begin{minipage}[t]{0.30\textwidth}
			\centering
			\CSampleHistogramTwenty
			\captionof{figure}{
				Histogram from a sample of $\CsixToneCount{n}=\CsixToneCount{20}$ observations.
			}
			\label{fig:c6-sample-histogram-20}
		\end{minipage}
		&
		\begin{minipage}[t]{0.30\textwidth}
			\centering
			\CSampleHistogramHundred
			\captionof{figure}{
				Histogram from a sample of $\CsixToneCount{n}=\CsixToneCount{100}$ observations.
			}
			\label{fig:c6-sample-histogram-100}
		\end{minipage}
		&
		\begin{minipage}[t]{0.30\textwidth}
			\centering
			\CSampleHistogramThousand
			\captionof{figure}{
				Histogram from a sample of $\CsixToneCount{n}=\CsixToneCount{1000}$ observations.
			}
			\label{fig:c6-sample-histogram-1000}
		\end{minipage}
	\end{tabular}
\end{center}

Histograms vary between samples. Larger samples generally give a more stable
approximation of the distribution's shape.

Compare the \CfiveTonePopulationMean{theoretical} \CfiveTonePopulationMean{mean} $\CfiveTonePopulationMean{\mu}$ and variance $\CfiveTonePopulationVariance{\sigma^2}$ with the sample
mean $\CsevenToneSampleMean{\bar{x}}$ and variance $\CsevenToneSampleVariance{s^2}$:

$$
\boxed{
	\begin{array}{\ConeToneVariable{c}|\ConeToneVariable{c}}
		\textbf{Theoretical model} & \textbf{Observed data} \\ \hline
		\text{Probability distribution}
		&
		\text{Observations }\ConeToneVariable{x_1},\ConeToneVariable{x_2},\ldots,\ConeToneVariable{x_n}
		\\[1mm]
		\text{Theoretical mean }\CfiveTonePopulationMean{\mu}
		&
		\text{Sample mean }\CsevenToneSampleMean{\bar{x}}
		\\[1mm]
		\text{Theoretical variance }\CfiveTonePopulationVariance{\sigma^2}
		&
		\text{Sample variance }\CsevenToneSampleVariance{s^2}
\end{array}}
$$

Larger samples generally give more stable estimates of these parameters:

$$
\CsevenToneSampleMean{\bar{x}}\approx\CfiveTonePopulationMean{\mu},
\qquad
\CsevenToneSampleVariance{s^2}\approx\CfiveTonePopulationVariance{\sigma^2}.
$$

The relationship is

$$
\boxed{
	\text{Real-world process}
	\;\longrightarrow\;
	\text{Observed data}
	\;\longrightarrow\;
	\text{Histogram and sample statistics}
	\;\longleftrightarrow\;
	\text{Probability distribution}
}
$$

\secInicio[sec:c6-intro-samples-histograms][sec:c6-introduction]

\subsection{Continuous distributions, theoretical and empirical probabilities}
\label{sec:c6-intro-distributions}

\begin{center}
	\CNormalPDFAndHistogram{0}{1}
	\captionof{figure}{
		Theoretical normal \CtwoToneDensity{probability} \CtwoToneDensity{density} \CtwoToneDensity{function} and a histogram
		from a sample generated from the same distribution. The smooth
		curve represents the theoretical model, while the bars represent
		the observed sample.
	}
	\label{fig:c6-normal-pdf-histogram}
\end{center}

For a continuous \ConeToneVariable{random} \ConeToneVariable{variable}, the theoretical probability of an interval
$[\ConeToneVariable{a},\ConeToneVariable{b}]$ is the area under its PDF:

$$
\ConeToneProbability{P}(\ConeToneVariable{a}\leq \ConeToneVariable{X}\leq \ConeToneVariable{b})
=
\int_a^b \CtwoToneDensity{f}(\ConeToneVariable{x})\,\ConeToneVariable{d}\ConeToneVariable{x}.
$$

If $\CsixToneCount{k}$ of $\CsixToneCount{n}$ sample observations fall within the interval, the empirical
probability is

$$
\CsixToneEmpirical{\widehat{P}}(\ConeToneVariable{a}\leq \ConeToneVariable{X}\leq \ConeToneVariable{b})
=
\frac{\CsixToneCount{k}}{\CsixToneCount{n}},
\qquad
\CsixToneCount{k}=\#\{\ConeToneVariable{x_i}:\ConeToneVariable{a}\leq \ConeToneVariable{x_i}\leq \ConeToneVariable{b}\}.
}$$

Compare the two methods:

$$
\boxed{
	\begin{array}{\ConeToneVariable{c}|\ConeToneVariable{c}}
		\textbf{Theoretical probability} & \textbf{Empirical probability} \\ \hline
		\text{Based on the probability distribution}
		&
		\text{Based on a generated sample}
		\\[1mm]
		\text{Area under the PDF}
		&
		\text{Relative frequency}
		\\[1mm]
		\displaystyle
		\ConeToneProbability{P}(\ConeToneVariable{a}\leq \ConeToneVariable{X}\leq \ConeToneVariable{b})
		=
		\int_a^b \CtwoToneDensity{f}(\ConeToneVariable{x})\,\ConeToneVariable{d}\ConeToneVariable{x}
		&
		\displaystyle
		\CsixToneEmpirical{\widehat{P}}(\ConeToneVariable{a}\leq \ConeToneVariable{X}\leq \ConeToneVariable{b})
		=
		\frac{\CsixToneCount{k}}{\CsixToneCount{n}}
\end{array}}
$$

The graphs illustrate each method:
\vspace{-4mm}
\begin{center}
	\begin{tabular}{cc}
		\begin{minipage}[t]{0.48\textwidth}
			\centering
			\CNormalPDFGraph{0}{1}
			\captionof{figure}{
				Theoretical \CtwoToneDensity{probability} \CtwoToneDensity{density} \CtwoToneDensity{function} of a normal
				distribution. The area under the curve over an interval
				represents the theoretical probability of that interval.
			}
			\label{fig:c6-normal-pdf}
		\end{minipage}
		&
		\begin{minipage}[t]{0.48\textwidth}
			\centering
			\CSampleHistogramNormal{0}{1}
			\captionof{figure}{
				Histogram of a finite sample generated from the same
				distribution. The bars show the observed frequencies
				within the histogram bins.
			}
			\label{fig:c6-normal-histogram}
		\end{minipage}
	\end{tabular}
\end{center}

In a histogram, observations within the selected interval determine $\CsixToneCount{k}$.
For every distribution, calculate \CsixToneEmpirical{empirical} \CsixToneEmpirical{probability} using

$$
\CsixToneEmpirical{\widehat{P}}
=
\frac{\text{number of observations in the interval}}
{\text{total number of observations}}.
$$

Theoretical probability depends on the distribution's PDF:

$$
\boxed{
	\text{PDF area}
	\quad\longleftrightarrow\quad
	\text{histogram relative frequency}
}
$$

These quantities generally become closer as \CsixToneCount{sample} \CsixToneCount{size} increases, though
random sampling causes differences.

The difference between the empirical and theoretical probabilities is

$$
\text{Difference}
=
\CsixToneEmpirical{\widehat{P}}(\ConeToneVariable{a}\leq \ConeToneVariable{X}\leq \ConeToneVariable{b})
-
\ConeToneProbability{P}(\ConeToneVariable{a}\leq \ConeToneVariable{X}\leq \ConeToneVariable{b}).
$$

The \CsixToneAbsoluteError{absolute} \CsixToneAbsoluteError{difference} is

$$
\text{Absolute difference}
=
\left|
\CsixToneEmpirical{\widehat{P}}(\ConeToneVariable{a}\leq \ConeToneVariable{X}\leq \ConeToneVariable{b})
-
\ConeToneProbability{P}(\ConeToneVariable{a}\leq \ConeToneVariable{X}\leq \ConeToneVariable{b})
\right|.
$$

When the theoretical probability is nonzero, the \CsixToneRelativeError{percentage} \CsixToneRelativeError{error} is

$$
\text{Percentage error}
=
\frac{
	\left|
	\CsixToneEmpirical{\widehat{P}}(\ConeToneVariable{a}\leq \ConeToneVariable{X}\leq \ConeToneVariable{b})
	-
	\ConeToneProbability{P}(\ConeToneVariable{a}\leq \ConeToneVariable{X}\leq \ConeToneVariable{b})
	\right|
}{
	\ConeToneProbability{P}(\ConeToneVariable{a}\leq \ConeToneVariable{X}\leq \ConeToneVariable{b})
}
\times100\%.
$$

\secInicio[sec:c6-intro-distributions][sec:c6-introduction]

\subsection{Sample mean and variance from generated data}
\label{sec:c6-intro-sample-mean-variance}

Use the following notation for theoretical parameters and sample statistics:

$$
\boxed{
	\begin{array}{\ConeToneVariable{c}|\ConeToneVariable{c}}
		\textbf{Theoretical model} & \textbf{Generated sample} \\ \hline
		\text{Probability distribution} & \text{Finite dataset}\\
		\text{Mean } \CfiveTonePopulationMean{\mu} & \text{Sample mean } \CsevenToneSampleMean{\bar{x}}\\
		\text{Variance } \CfiveTonePopulationVariance{\sigma^2} & \text{Sample variance }\CsevenToneSampleVariance{s^2}\\
		\text{Standard deviation } \CfiveTonePopulationSD{\sigma} & \text{Sample standard deviation }\CsevenToneSampleSD{s}
\end{array}}
$$
For ungrouped observations $\ConeToneVariable{x_1},\ConeToneVariable{x_2},\ldots,\ConeToneVariable{x_n}$, the \CsevenToneSampleMean{sample} \CsevenToneSampleMean{mean}, variance,
and standard deviation are

$$
\CsevenToneSampleMean{\bar{x}}
=
\frac{\ConeToneVariable{1}}{\CsixToneCount{n}}\sum_{i=1}^{n}\ConeToneVariable{x_i},
\qquad
\CsevenToneSampleVariance{s^2}
=
\frac{\ConeToneVariable{1}}{\CsixToneCount{n}\CsixToneCount{-1}}\sum_{i=1}^{n}(\ConeToneVariable{x_i}-\CsevenToneSampleMean{\bar{x}})^2,
\qquad
\CsevenToneSampleSD{s}=\sqrt{\CsevenToneSampleVariance{s^2}}.
$$

The mean measures the data's center; the variance measures spread around
$\CsevenToneSampleMean{\bar{x}}$.

When observations are grouped into intervals, let $\CsixToneCount{f_i}$ be the frequency of
class $i$ and $\ConeToneVariable{m_i}$ its midpoint:

$$
\ConeToneVariable{m_i}=\frac{\ConeToneVariable{a_i}+\ConeToneVariable{b_i}}{2},
\qquad
\CsixToneCount{n}=\sum_i \CsixToneCount{f_i}.
$$

The grouped-data mean and variance are

$$
\CsevenToneSampleMean{\bar{x}}
=
\frac{\sum_i \CsixToneCount{f_i} \ConeToneVariable{m_i}}{\CsixToneCount{n}},
\qquad
\CsevenToneSampleVariance{s^2}
=
\frac{\sum_i \CsixToneCount{f_i}(\ConeToneVariable{m_i}-\CsevenToneSampleMean{\bar{x}})^2}{\CsixToneCount{n}\CsixToneCount{-1}}.
$$

Grouped calculations are approximate because all observations within an
interval are represented by its midpoint.

$$
\boxed{
	\begin{array}{\ConeToneVariable{c}|\ConeToneVariable{c}}
		\textbf{Ungrouped data} & \textbf{Grouped data}\\ \hline
		\ConeToneVariable{x_1},\ConeToneVariable{x_2},\ldots,\ConeToneVariable{x_n} & \text{Intervals and frequencies}\\
		\text{Exact observations} & \text{Class midpoints }\ConeToneVariable{m_i}\\
		\displaystyle
		\CsevenToneSampleMean{\bar{x}}=\frac{\ConeToneVariable{1}}{\CsixToneCount{n}}\sum_i \ConeToneVariable{x_i}
		&
		\displaystyle
		\CsevenToneSampleMean{\bar{x}}=\frac{\sum_i \CsixToneCount{f_i} \ConeToneVariable{m_i}}{\CsixToneCount{n}}
		\\[\ConeToneVariable{3}mm]
		\displaystyle
		\CsevenToneSampleVariance{s^2}=\frac{\ConeToneVariable{1}}{\CsixToneCount{n}-1}\sum_i(\ConeToneVariable{x_i}-\CsevenToneSampleMean{\bar{x}})^2
		&
		\displaystyle
		\CsevenToneSampleVariance{s^2}=\frac{\sum_i \CsixToneCount{f_i}(\ConeToneVariable{m_i}-\CsevenToneSampleMean{\bar{x}})^2}{\CsixToneCount{n}-1}
\end{array}}
$$

For a finite random sample generated from a theoretical distribution,

$$
\CsevenToneSampleMean{\bar{x}}\approx\CfiveTonePopulationMean{\mu},
\qquad
\CsevenToneSampleVariance{s^2}\approx\CfiveTonePopulationVariance{\sigma^2}.
$$

Estimates vary between samples and generally become more stable as sample
size increases.

\secInicio[sec:c6-intro-sample-mean-variance][sec:c6-introduction]

% ==========================================================================
% C6-ONLY PRACTICE
% ==========================================================================


\Needspace{8\baselineskip}
\subsection*{Notation reference}
\phantomsection\label{tone:C6C7:reference}
\COneToneDetailedReference{C6C7}
\section{C6 Practice: Empirical and Theoretical Probabilities}
\label{sec:c6-only-practice}

\SectionProblemsTableOfContents{
	\SectionProblemLink{sec:c6-only-raw}{Monthly household electricity use (raw data)} \\
	\SectionProblemLink{sec:c6-only-frequency}{Waiting times (frequency data)} \\
	\SectionProblemLink{sec:c6-only-grouped}{Package masses (grouped data)} \\
	\SectionProblemLink{sec:c6-only-histogram}{Cafeteria waiting times (histogram data)}
}\vspace{5mm}

The following problems focus only on estimating interval probabilities from
samples. The larger samples make the relative-frequency estimates more stable,
while the raw, frequency, and grouped formats provide different ways of
organizing the observations.

\subsection{Monthly household electricity use (raw data)}
\label{sec:c6-only-raw}

A utility company generates the following sample of monthly household
electricity use (in hundreds of kWh):

\[
\begin{array}{cccccccccc}
	3.8 & 4.1 & 4.5 & 5.2 & 5.7 & 6.0 & 6.3 & 6.8 & 7.1 & 7.5 \\
	4.0 & 4.4 & 4.9 & 5.4 & 5.9 & 6.1 & 6.5 & 6.9 & 7.3 & 8.0
\end{array}
\]

For the theoretical model,
$P(5\leq X\leq7)=0.62$.

\QuestionsTasks[sec:c6-only-practice]{
	{6}/{subsec:c6-only-raw-solution}/
	{Estimate $P(5\leq X\leq7)$ from the sample. Then calculate the
	absolute error and the percentage error relative to the theoretical
	probability.
	}
}

\subsubsection*{C6 Solution}
\phantomsection
\label{subsec:c6-only-raw-solution}

There are $\CsixToneCount{n}=\CsixToneCount{20}$ observations. The values in the interval from $5$ to $7$,
inclusive, are

\[
5.2,\ 5.7,\ 6.0,\ 6.3,\ 6.8,\ 5.4,\ 5.9,\ 6.1,\ 6.5,\ 6.9,
\]

so $\CsixToneCount{k}=\CsixToneCount{10}$. Therefore,

\[
\CsixToneEmpirical{\widehat P}(5\leq \ConeToneVariable{X}\leq7)=\frac{\CsixToneCount{k}}{\CsixToneCount{n}}=\frac{\CsixToneCount{\ConeToneVariable{10}}}{\CsixToneCount{\CsixToneCount{20}}}=\CsixToneEmpirical{0.50}.
\]

The \CsixToneAbsoluteError{absolute} \CsixToneAbsoluteError{error} is

\[
\CsixToneAbsoluteError{E}=|\CsixToneEmpirical{0.50}-\ConeToneProbability{0.62}|=\CsixToneAbsoluteError{0.12},
\]

and the \CsixToneRelativeError{percentage} \CsixToneRelativeError{error} is

\[
\frac{|\ConeToneVariable{0.50}\ConeToneVariable{-0.62}|}{0.62}\times100\%
=\frac{\ConeToneVariable{0.12}}{0.62}\times100\%
\approx19.35\%.
\]

Thus the \CsixToneEmpirical{empirical} \CsixToneEmpirical{probability} is $\ConeToneProbability{0.50}$, which is $0.12$ below the
theoretical probability.

\secInicio[sec:c6-only-raw][sec:c6-only-practice]

\subsection{Waiting times (frequency data)}
\label{sec:c6-only-frequency}

A clinic records the waiting times of $120$ patients in the following
frequency table:

\[
\begin{array}{c|cccccc}
	\text{Waiting time (min)} & 5 & 10 & 15 & 20 & 25 & 30 \\ \hline
	\text{Frequency} & 8 & 20 & 34 & 30 & 18 & 10
\end{array}
\]

For the theoretical model,
$P(10\leq X\leq20)=0.72$.

\QuestionsTasks[sec:c6-only-practice]{
	{6}/{subsec:c6-only-frequency-solution}/
	{Estimate the probability that a patient waits from $10$ to $20$
	minutes, inclusive. Calculate its absolute error and compare the
	empirical estimate with the theoretical probability.
	}
}

\subsubsection*{C6 Solution}
\phantomsection
\label{subsec:c6-only-frequency-solution}

The total frequency is

\[
\CsixToneCount{n}=\CsixToneCount{8}+20+34+30+18+10=\CsixToneCount{120}.
\]

The waiting times $10$, $15$, and $20$ minutes are in the interval, so

\[
\CsixToneCount{k}=\CsixToneCount{20}+34+30=\CsixToneCount{84}.
\]

Hence

\[
\CsixToneEmpirical{\widehat P}(10\leq \ConeToneVariable{X}\leq20)=\frac{\ConeToneVariable{84}}{\CsixToneCount{120}}=\CsixToneEmpirical{0.70}.
\]

The \CsixToneAbsoluteError{absolute} \CsixToneAbsoluteError{error} is

\[
\CsixToneAbsoluteError{E}=|\CsixToneAbsoluteError{0.70}-0.72|=\CsixToneAbsoluteError{0.02}.
\]

Thus the \CsixToneEmpirical{empirical} \CsixToneEmpirical{probability} is $\ConeToneProbability{0.70}$, which is $0.02$ below the
theoretical probability of $\ConeToneProbability{0.72}$.

\secInicio[sec:c6-only-frequency][sec:c6-only-practice]

\subsection{Package masses (grouped data)}
\label{sec:c6-only-grouped}

A distribution center groups the masses of $200$ packages as follows:

\[
\begin{array}{c|ccccc}
	\text{Package mass (kg)}
	& 0\text{--}2 & 2\text{--}4 & 4\text{--}6 & 6\text{--}8 & 8\text{--}10 \\ \hline
	\text{Frequency} & 18 & 46 & 72 & 44 & 20
\end{array}
\]

The classes are left-inclusive and right-exclusive. For the theoretical
model, $P(2\leq X<8)=0.83$.

\QuestionsTasks[sec:c6-only-practice]{
	{6}/{subsec:c6-only-grouped-solution}/
	{Estimate $P(2\leq X<8)$ from the grouped sample. Calculate the
	absolute error and the percentage error relative to the theoretical
	probability.
	}
}

\subsubsection*{C6 Solution}
\phantomsection
\label{subsec:c6-only-grouped-solution}

The total frequency is

\[
\CsixToneCount{n}=\CsixToneCount{18}+46+72+44+20=\CsixToneCount{200}.
\]

The interval $2\leq \ConeToneVariable{X}<8$ contains the three classes $2$--$4$, $4$--$6$,
and $6$--$8$. Their combined frequency is

\[
\CsixToneCount{k}=\CsixToneCount{46}+72+44=\CsixToneCount{162}.
\]

Therefore,

\[
\CsixToneEmpirical{\widehat P}(2\leq \ConeToneVariable{X}<8)=\frac{\ConeToneVariable{162}}{\CsixToneCount{200}}=\CsixToneEmpirical{0.81}.
\]

The \CsixToneAbsoluteError{absolute} \CsixToneAbsoluteError{error} is

\[
\CsixToneAbsoluteError{E}=|\CsixToneAbsoluteError{0.81}-0.83|=\CsixToneAbsoluteError{0.02},
\]

and the \CsixToneRelativeError{percentage} \CsixToneRelativeError{error} is

\[
\frac{|\ConeToneVariable{0.81}\ConeToneVariable{-0.83}|}{0.83}\times100\%
=\frac{\ConeToneVariable{0.02}}{0.83}\times100\%
\approx2.41\%.
\]

Thus the empirical estimate is $0.81$, which is close to, but slightly
below, the theoretical probability of $\ConeToneProbability{0.83}$.

\secInicio[sec:c6-only-grouped][sec:c6-only-practice]

% ==========================================================================
% C7-ONLY PRACTICE
% ==========================================================================

\subsection{Cafeteria waiting times (histogram data)}
\label{sec:c6-only-histogram}

A school cafeteria records a sample of student waiting times $x$, in minutes,
during lunch. The sample is presented only in the histogram below. Bar heights
represent frequencies, and all classes are left-inclusive and right-exclusive.

\HistogramGraph{(0,2),(2,4),(4,6),(6,8)}{4,8,6,2}
    {0.70\textwidth}{5cm}{green!25}{2}{2}

\begin{center}
    Waiting time $x$ (minutes); vertical axis: frequency (number of students).
\end{center}

For the theoretical model, $P(2\leq X<6)=0.75$.

\QuestionsTasks[sec:c6-only-practice]{
    {6}/{subsec:c6-only-histogram-solution}/
    {Read each class interval and frequency from the histogram, and determine
    the sample size. Calculate $\widehat P(2\leq X<6)$ and
    $\widehat P(X\geq6)$. Find the cumulative frequency and cumulative
    empirical probability below each upper class boundary. Compare
    $\widehat P(2\leq X<6)$ with the theoretical probability by calculating
    the absolute error and percentage error.
    }
}

\subsubsection*{C6 Solution}
\phantomsection
\label{subsec:c6-only-histogram-solution}

First, read the class boundaries on the horizontal axis and the bar heights
on the vertical axis. The frequency grid spacing is $2$:

\[
\begin{array}{\ConeToneVariable{c}|\ConeToneVariable{c}\ConeToneVariable{c}\ConeToneVariable{c}\ConeToneVariable{c}}
    \text{Waiting time (min)} & [0,2) & [2,4) & [4,6) & [6,8) \\ \hline
    \CsixToneCount{f_i} & 4 & 8 & 6 & 2
\end{array}
\]

The total frequency gives the \CsixToneCount{sample} \CsixToneCount{size}:

\[
\CsixToneCount{n}=\sum_i \CsixToneCount{f_i}=\CsixToneCount{4}+8+6+2=\CsixToneCount{20}.
\]

The interval $[2,6)$ contains two complete classes, so its frequency is
$\CsixToneCount{k}=\CsixToneCount{8}+6=\CsixToneCount{14}$. Therefore,

\[
\CsixToneEmpirical{\widehat P}(2\leq \ConeToneVariable{X}<6)=\frac{\CsixToneCount{k}}{\CsixToneCount{n}}=\frac{\ConeToneVariable{14}}{\CsixToneCount{20}}=\CsixToneEmpirical{0.70}.
\]

Only the last class contains sampled waiting times at least $6$ minutes:

\[
\CsixToneEmpirical{\widehat P}(\ConeToneVariable{X}\geq6)=\frac{\ConeToneVariable{2}}{\CsixToneCount{20}}=\CsixToneEmpirical{0.10}.
\]

For each upper boundary $\ConeToneVariable{b}$, add all frequencies below $\ConeToneVariable{b}$, then divide
by $20$ to calculate the cumulative \CsixToneEmpirical{empirical} \CsixToneEmpirical{probability}:

\[
\begin{array}{\ConeToneVariable{c}|\ConeToneVariable{c}|\ConeToneVariable{c}}
    \ConeToneVariable{b}\text{ (min)} & \text{Cumulative frequency} & \CsixToneEmpirical{\widehat P}(\ConeToneVariable{X}<\ConeToneVariable{b}) \\ \hline
    2 & 4 & 4/20=\ConeToneVariable{0.20} \\
    4 & 4+8=\ConeToneVariable{12} & 12/20=\ConeToneVariable{0.60} \\
    6 & 4+8+6=\ConeToneVariable{18} & 18/20=\ConeToneVariable{0.90} \\
    8 & 4+8+6+2=\ConeToneVariable{20} & 20/20=\ConeToneVariable{1}
\end{array}
\]

Finally, compare the \CsixToneEmpirical{empirical} \CsixToneEmpirical{probability} $\ConeToneProbability{0.70}$ with the theoretical
probability $\ConeToneProbability{0.75}$. The \CsixToneAbsoluteError{absolute} \CsixToneAbsoluteError{error} is

\[
\CsixToneAbsoluteError{E}=|\CsixToneAbsoluteError{0.70}-0.75|=\CsixToneAbsoluteError{0.05},
\]

and the \CsixToneRelativeError{percentage} \CsixToneRelativeError{error} is

\[
\frac{|\ConeToneVariable{0.70}\ConeToneVariable{-0.75}|}{0.75}\times100\%
=\frac{\ConeToneVariable{0.05}}{0.75}\times100\%\approx6.67\%.
\]

Thus the \CsixToneEmpirical{empirical} \CsixToneEmpirical{probability} is $\ConeToneProbability{0.05}$ below the theoretical probability.

\secInicio[sec:c6-only-histogram][sec:c6-only-practice]

\section{C7 Practice: Sample Mean, Variance, and Standard Deviation}
\label{sec:c7-only-practice}

\SectionProblemsTableOfContents{
	\SectionProblemLink{sec:c7-only-raw}{Machine calibration errors (raw data)} \\
	\SectionProblemLink{sec:c7-only-frequency}{Books borrowed (frequency data)} \\
	\SectionProblemLink{sec:c7-only-grouped}{Commute times (grouped data)} \\
	\SectionProblemLink{sec:c7-only-histogram}{Parcel masses (histogram data)}
}\vspace{5mm}

The samples in this section are deliberately small so that the sample mean,
sample variance, and sample standard deviation can be calculated clearly by
hand. Recall that the sample-variance denominator is $n-1$.

\subsection{Machine calibration errors (raw data)}
\label{sec:c7-only-raw}

The calibration errors (in tenths of a millimeter) from five trials are

\[
4,\quad5,\quad6,\quad7,\quad8.
\]

The theoretical model has $\mu=6$ and $\sigma^2=2.25$.

\QuestionsTasks[sec:c7-only-practice]{
	{7}/{subsec:c7-only-raw-solution}/
	{Calculate the sample mean, sample variance, and sample standard
	deviation. Compare the sample mean and variance with the theoretical
	values.
	}
}

\subsubsection*{C7 Solution}
\phantomsection
\label{subsec:c7-only-raw-solution}

There are $\CsixToneCount{n}=\CsixToneCount{5}$ observations. First, calculate the \CsevenToneSampleMean{sample} \CsevenToneSampleMean{mean}:

\[
\CsevenToneSampleMean{\bar{x}}
=\frac{\ConeToneVariable{4}+\ConeToneVariable{5}+\ConeToneVariable{6}+\ConeToneVariable{7}+\ConeToneVariable{8}}{\CsixToneCount{5}}
=\frac{\ConeToneVariable{30}}{\CsixToneCount{5}}
=\CsevenToneSampleMean{6}.
\]

Next, calculate the squared deviations from $\CsevenToneSampleMean{\bar{x}}=\CsevenToneSampleMean{6}$:

\[
\begin{array}{\ConeToneVariable{c}|\ConeToneVariable{c}\ConeToneVariable{c}\ConeToneVariable{c}\ConeToneVariable{c}\ConeToneVariable{c}}
	\ConeToneVariable{x_i} & 4 & 5 & 6 & 7 & 8 \\ \hline
	\ConeToneVariable{x_i}-\CsevenToneSampleMean{\bar{x}} & -2 & -1 & 0 & 1 & 2 \\
	(\ConeToneVariable{x_i}-\CsevenToneSampleMean{\bar{x}})^2 & 4 & 1 & 0 & 1 & 4
\end{array}
\]

Their sum is $4+1+0+1+4=10$. Therefore, the \CsevenToneSampleVariance{sample} \CsevenToneSampleVariance{variance} is

\[
\CsevenToneSampleVariance{s^2}
=\frac{\sum_{i=1}^{5}(\ConeToneVariable{x_i}-\CsevenToneSampleMean{\bar{x}})^2}{\CsixToneCount{n}\CsixToneCount{-1}}
=\frac{\ConeToneVariable{10}}{\CsixToneCount{5}\CsixToneCount{-1}}
=\CsevenToneSampleVariance{2.50}.
\]

The \CsevenToneSampleSD{sample} \CsevenToneSampleSD{standard} \CsevenToneSampleSD{deviation} is

\[
\CsevenToneSampleSD{s}=\sqrt{\CsevenToneSampleVariance{s^2}}=\sqrt{\CsevenToneSampleVariance{2.50}}\approx\CsevenToneSampleSD{1.58}.
\]

Thus,

\[
\boxed{\CsevenToneSampleMean{\bar{x}}=\CsevenToneSampleMean{6}=\CfiveTonePopulationMean{\mu},\qquad
\CsevenToneSampleVariance{s^2}=\CsevenToneSampleVariance{2.50}\approx\CfiveTonePopulationVariance{\sigma^2}=\CsevenToneSampleVariance{2.25},\qquad \CsevenToneSampleSD{s}\approx\CsevenToneSampleSD{1.58}.}
\]

\secInicio[sec:c7-only-raw][sec:c7-only-practice]

\subsection{Books borrowed (frequency data)}
\label{sec:c7-only-frequency}

The numbers of books borrowed by four students are summarized below:

\[
\begin{array}{c|ccc}
	x_i & 2 & 4 & 6 \\ \hline
	f_i & 1 & 2 & 1
\end{array}
\]

The theoretical model has $\mu=4$ and $\sigma^2=3$.

\QuestionsTasks[sec:c7-only-practice]{
	{7}/{subsec:c7-only-frequency-solution}/
	{Calculate the sample mean, sample variance, and sample standard
	deviation from the frequency table. Compare the sample mean and
	variance with the theoretical values.
	}
}

\subsubsection*{C7 Solution}
\phantomsection
\label{subsec:c7-only-frequency-solution}

The total frequency, and hence the \CsixToneCount{sample} \CsixToneCount{size}, is

\[
\CsixToneCount{n}=\CsixToneCount{1}+2+1=\CsixToneCount{4}.
\]

The \CsevenToneSampleMean{sample} \CsevenToneSampleMean{mean} is

\[
\CsevenToneSampleMean{\bar{x}}
=\frac{\sum_i \CsixToneCount{f_i} \ConeToneVariable{x_i}}{\CsixToneCount{n}}
=\frac{\ConeToneVariable{1}(\ConeToneVariable{2})+\ConeToneVariable{2}(\ConeToneVariable{4})+\ConeToneVariable{1}(\ConeToneVariable{6})}{\CsixToneCount{4}}
=\frac{\ConeToneVariable{16}}{\CsixToneCount{4}}
=\CsevenToneSampleMean{4}.
\]

Using $\CsevenToneSampleMean{\bar{x}}=\CsevenToneSampleMean{4}$, calculate each weighted squared deviation:

\[
\begin{array}{\ConeToneVariable{c}|\ConeToneVariable{c}\ConeToneVariable{c}\ConeToneVariable{c}}
	\ConeToneVariable{x_i} & 2 & 4 & 6 \\ \hline
	\CsixToneCount{f_i} & 1 & 2 & 1 \\
	(\ConeToneVariable{x_i}-\CsevenToneSampleMean{\bar{x}})^2 & 4 & 0 & 4 \\
	\CsixToneCount{f_i}(\ConeToneVariable{x_i}-\CsevenToneSampleMean{\bar{x}})^2 & 4 & 0 & 4
\end{array}
\]

The sum of the weighted squared deviations is $4+0+4=8$. Thus,

\[
\CsevenToneSampleVariance{s^2}
=\frac{\sum_i \CsixToneCount{f_i}(\ConeToneVariable{x_i}-\CsevenToneSampleMean{\bar{x}})^2}{\CsixToneCount{n}\CsixToneCount{-1}}
=\frac{\ConeToneVariable{8}}{\CsixToneCount{4}\CsixToneCount{-1}}
=\frac83
\approx\CsevenToneSampleVariance{2.67}.
\]

The \CsevenToneSampleSD{sample} \CsevenToneSampleSD{standard} \CsevenToneSampleSD{deviation} is

\[
\CsevenToneSampleSD{s}=\sqrt{\CsevenToneSampleVariance{s^2}}=\sqrt{\frac83}\approx\CsevenToneSampleSD{1.63}.
\]

Therefore,

\[
\boxed{\CsevenToneSampleMean{\bar{x}}=\CsevenToneSampleMean{4}=\CfiveTonePopulationMean{\mu},\qquad
\CsevenToneSampleVariance{s^2}\approx\CsevenToneSampleVariance{2.67}\approx\CfiveTonePopulationVariance{\sigma^2}=\CsevenToneSampleVariance{3},\qquad \CsevenToneSampleSD{s}\approx\CsevenToneSampleSD{1.63}.}
\]

\secInicio[sec:c7-only-frequency][sec:c7-only-practice]

\subsection{Commute times (grouped data)}
\label{sec:c7-only-grouped}

The commute times of five employees are grouped into three classes:

\[
\begin{array}{c|ccc}
	\text{Commute time (min)} & 0\text{--}10 & 10\text{--}20 & 20\text{--}30 \\ \hline
	\text{Frequency} & 1 & 3 & 1
\end{array}
\]

The theoretical model has $\mu=15$ minutes and
$\sigma^2=48$ minutes$^2$.

\QuestionsTasks[sec:c7-only-practice]{
	{7}/{subsec:c7-only-grouped-solution}/
	{Use class midpoints to estimate the sample mean, sample variance, and
	sample standard deviation. Compare the estimated mean and variance
	with the theoretical values.
	}
}

\subsubsection*{C7 Solution}
\phantomsection
\label{subsec:c7-only-grouped-solution}

First, calculate each class midpoint:

\[
\ConeToneVariable{m_1}=\frac{\ConeToneVariable{0}+\ConeToneVariable{10}}{2}=\ConeToneVariable{5},\qquad
\ConeToneVariable{m_2}=\frac{\ConeToneVariable{10}+\ConeToneVariable{20}}{2}=\ConeToneVariable{15},\qquad
\ConeToneVariable{m_3}=\frac{\ConeToneVariable{20}+\ConeToneVariable{30}}{2}=\ConeToneVariable{25}.
\]

The total frequency is

\[
\CsixToneCount{n}=\CsixToneCount{1}+3+1=\CsixToneCount{5}.
\]

Using the midpoints as representative values, the estimated \CsevenToneSampleMean{sample} \CsevenToneSampleMean{mean} is

\[
\CsevenToneSampleMean{\bar{x}}
=\frac{\sum_i \CsixToneCount{f_i} \ConeToneVariable{m_i}}{\CsixToneCount{n}}
=\frac{\ConeToneVariable{1}(\ConeToneVariable{5})+\ConeToneVariable{3}(\ConeToneVariable{15})+\ConeToneVariable{1}(\ConeToneVariable{25})}{\CsixToneCount{5}}
=\frac{\ConeToneVariable{75}}{\CsixToneCount{5}}
=\CsevenToneSampleMean{15}\text{ min}.
\]

Next, calculate the weighted squared deviations:

\[
\begin{array}{\ConeToneVariable{c}|\ConeToneVariable{c}\ConeToneVariable{c}\ConeToneVariable{c}}
	\ConeToneVariable{m_i} & 5 & 15 & 25 \\ \hline
	\CsixToneCount{f_i} & 1 & 3 & 1 \\
	(\ConeToneVariable{m_i}-\CsevenToneSampleMean{\bar{x}})^2 & 100 & 0 & 100 \\
	\CsixToneCount{f_i}(\ConeToneVariable{m_i}-\CsevenToneSampleMean{\bar{x}})^2 & 100 & 0 & 100
\end{array}
\]

Their sum is $100+0+100=200$. The estimated \CsevenToneSampleVariance{sample} \CsevenToneSampleVariance{variance} is

\[
\CsevenToneSampleVariance{s^2}
=\frac{\sum_i \CsixToneCount{f_i}(\ConeToneVariable{m_i}-\CsevenToneSampleMean{\bar{x}})^2}{\CsixToneCount{n}\CsixToneCount{-1}}
=\frac{\ConeToneVariable{200}}{\CsixToneCount{5}\CsixToneCount{-1}}
=\CsevenToneSampleVariance{50}\text{ min}^2.
\]

The estimated \CsevenToneSampleSD{sample} \CsevenToneSampleSD{standard} \CsevenToneSampleSD{deviation} is

\[
\CsevenToneSampleSD{s}=\sqrt{\CsevenToneSampleVariance{s^2}}=\sqrt{\CsevenToneSampleVariance{50}}\approx\CsevenToneSampleSD{7.07}\text{ min}.
\]

Therefore,

\[
\boxed{\CsevenToneSampleMean{\bar{x}}=\CsevenToneSampleMean{15}\text{ min}=\CfiveTonePopulationMean{\mu},\qquad
\CsevenToneSampleVariance{s^2}=\CsevenToneSampleVariance{50}\text{ min}^2\approx\CfiveTonePopulationVariance{\sigma^2}=\CsevenToneSampleVariance{48}\text{ min}^2,\qquad
\CsevenToneSampleSD{s}\approx\CsevenToneSampleSD{7.07}\text{ min}.}
\]

These results are estimates because all observations in a class are
represented by the class midpoint.

\secInicio[sec:c7-only-grouped][sec:c7-only-practice]

\subsection{Parcel masses (histogram data)}
\label{sec:c7-only-histogram}

A courier weighs a small sample of parcels. Their masses $x$, in kilograms,
are presented only in the histogram below. Bar heights represent frequencies,
and all classes are left-inclusive and right-exclusive.

\HistogramGraph{(0,2),(2,4),(4,6)}{2,4,2}
    {0.65\textwidth}{5cm}{orange!30}{2}{1}

\begin{center}
    Parcel mass $x$ (kg); vertical axis: frequency (number of parcels).
\end{center}

\QuestionsTasks[sec:c7-only-practice]{
    {7}/{subsec:c7-only-histogram-solution}/
    {Read the class intervals and frequencies from the histogram, and
    determine the sample size and class midpoints. Use the midpoints to
    estimate the sample mean, sample variance (with denominator $n-1$),
    and sample standard deviation. Explain why these statistics are estimates.
    }
}

\subsubsection*{C7 Solution}
\phantomsection
\label{subsec:c7-only-histogram-solution}

First, read the class boundaries from the horizontal axis and the frequencies
from the bar heights. The frequency grid spacing is $1$:

\[
\begin{array}{\ConeToneVariable{c}|\ConeToneVariable{c}\ConeToneVariable{c}\ConeToneVariable{c}}
    \text{Parcel mass (kg)} & [0,2) & [2,4) & [4,6) \\ \hline
    \CsixToneCount{f_i} & 2 & 4 & 2
\end{array}
\]

The \CsixToneCount{sample} \CsixToneCount{size} is the total frequency:

\[
\CsixToneCount{n}=\CsixToneCount{2}+4+2=\CsixToneCount{8}.
\]

Next, calculate each class midpoint using $\ConeToneVariable{m_i}=(\ConeToneVariable{a_i}+\ConeToneVariable{b_i})/\ConeToneVariable{2}$:

\[
\ConeToneVariable{m_1}=\frac{\ConeToneVariable{0}+\ConeToneVariable{2}}{2}=\ConeToneVariable{1},\qquad
\ConeToneVariable{m_2}=\frac{\ConeToneVariable{2}+\ConeToneVariable{4}}{2}=\ConeToneVariable{3},\qquad
\ConeToneVariable{m_3}=\frac{\ConeToneVariable{4}+\ConeToneVariable{6}}{2}=\ConeToneVariable{5}.
\]

Using these representative values, the estimated \CsevenToneSampleMean{sample} \CsevenToneSampleMean{mean} is

\[
\CsevenToneSampleMean{\bar{x}}=\frac{\sum_i \CsixToneCount{f_i} \ConeToneVariable{m_i}}{\CsixToneCount{n}}
=\frac{\ConeToneVariable{2}(\ConeToneVariable{1})+\ConeToneVariable{4}(\ConeToneVariable{3})+\ConeToneVariable{2}(\ConeToneVariable{5})}{\CsixToneCount{8}}
=\frac{\ConeToneVariable{2}+\ConeToneVariable{12}+\ConeToneVariable{10}}{\CsixToneCount{8}}
=\frac{\ConeToneVariable{24}}{\CsixToneCount{8}}=\CsevenToneSampleMean{3}\text{ kg}.
\]

Using $\CsevenToneSampleMean{\bar{x}}=\CsevenToneSampleMean{3}$, calculate the weighted squared deviations:

\[
\begin{array}{\ConeToneVariable{c}|\ConeToneVariable{c}\ConeToneVariable{c}\ConeToneVariable{c}}
    \ConeToneVariable{m_i} & 1 & 3 & 5 \\ \hline
    \CsixToneCount{f_i} & 2 & 4 & 2 \\
    \ConeToneVariable{m_i}-\CsevenToneSampleMean{\bar{x}} & -2 & 0 & 2 \\
    (\ConeToneVariable{m_i}-\CsevenToneSampleMean{\bar{x}})^2 & 4 & 0 & 4 \\
    \CsixToneCount{f_i}(\ConeToneVariable{m_i}-\CsevenToneSampleMean{\bar{x}})^2 & 8 & 0 & 8
\end{array}
\]

Their sum is $8+0+8=16$. The estimated \CsevenToneSampleVariance{sample} \CsevenToneSampleVariance{variance} uses $\CsixToneCount{n}\CsevenToneSampleVariance{-1}=\CsixToneCount{7}$:

\[
\CsevenToneSampleVariance{s^2}=\frac{\sum_i \CsixToneCount{f_i}(\ConeToneVariable{m_i}-\CsevenToneSampleMean{\bar{x}})^2}{\CsixToneCount{n}\CsixToneCount{-1}}
=\frac{\ConeToneVariable{16}}{\CsixToneCount{8}\CsixToneCount{-1}}=\frac{\ConeToneVariable{16}}{\CsixToneCount{7}}\approx\CsevenToneSampleVariance{2.29}\text{ kg}^2.
\]

Take the square root of the unrounded variance to estimate the standard deviation:

\[
\CsevenToneSampleSD{s}=\sqrt{\CsevenToneSampleVariance{s^2}}=\sqrt{\frac{\ConeToneVariable{16}}{7}}\approx\CsevenToneSampleSD{1.51}\text{ kg}.
\]

Therefore, the midpoint estimates are

\[
\boxed{\CsevenToneSampleMean{\bar{x}}=\CsevenToneSampleMean{3}\text{ kg},\qquad
\CsevenToneSampleVariance{s^2}\approx\CsevenToneSampleVariance{2.29}\text{ kg}^2,\qquad \CsevenToneSampleSD{s}\approx\CsevenToneSampleSD{1.51}\text{ kg}.}
\]

These statistics are estimates because the histogram gives grouped masses,
and every observation in a class is represented by its midpoint rather
than its individual mass.

\secInicio[sec:c7-only-histogram][sec:c7-only-practice]

\section{Samples, Empirical Probabilities, and Sample Statistics}
\label{sec:c6-specific-statistics}

\SectionProblemsTableOfContents{
	\SectionProblemLink{sec:c6c7-data-raw}{Investment returns (raw data)} \\
	\SectionProblemLink{sec:c6-probability-frequency}{Daily energy consumption (frequency data)} \\
	\SectionProblemLink{sec:c6-probability-grouped}{Delivery times (grouped data)} \\
	\SectionProblemLink{sec:c6c7-histogram}{Library visit durations (histogram data)}
}\vspace{5mm}

The following examples use small samples so that each step of the calculations can be shown clearly. In real-world applications, sample sizes are typically much larger to obtain more reliable estimates.

\subsection{Investment returns (raw data)}
\label{sec:c6c7-data-raw}

A financial analyst records these returns (in percent) over six days:

\[
2,\quad 2,\quad 3,\quad 3,\quad 4,\quad 4.
\]

The theoretical model has $\mu=3$, $\sigma^2=1$, and
$P(2\leq X\leq4)=0.80$.

\QuestionsTasks[sec:c6-specific-statistics]{
	{6}/{subsec:c6-probability-raw-b}/
	{Calculate the empirical probability of a return between $2\%$ and
	$4\%$, inclusive, and its absolute error.
	},
	{7}/{subsec:c7-meanvar-raw-b}/
	{Calculate the sample mean and sample variance, and compare them with
	the theoretical values.
	}
}

\subsubsection*{C6 Solution}
\phantomsection
\label{subsec:c6-probability-raw-b}

All six observations are in the interval, so $\CsixToneCount{n}=\CsixToneCount{6}$ and $\CsixToneCount{k}=\CsixToneCount{6}$. Hence

\[
\CsixToneEmpirical{\widehat P}(2\leq \ConeToneVariable{X}\leq4)=\frac{\CsixToneCount{k}}{\CsixToneCount{n}}=\frac{\ConeToneVariable{6}}{\CsixToneCount{6}}=\CsixToneEmpirical{1}.
\]

The \CsixToneAbsoluteError{absolute} \CsixToneAbsoluteError{error} is

\[
\CsixToneAbsoluteError{E}=|\CsixToneEmpirical{\widehat P}-\ConeToneProbability{P}|=|\CsixToneAbsoluteError{1}-0.80|=\CsixToneAbsoluteError{0.20}.
\]

Thus the \CsixToneEmpirical{empirical} \CsixToneEmpirical{probability} is $\ConeToneProbability{1}$, which is $0.20$ above the
theoretical probability.

\secInicio[sec:c6c7-data-raw][sec:c6-specific-statistics]

\subsubsection*{C7 Solution}
\phantomsection
\label{subsec:c7-meanvar-raw-b}

The \CsevenToneSampleMean{sample} \CsevenToneSampleMean{mean} is

\[
\CsevenToneSampleMean{\bar{x}}
=\frac{\ConeToneVariable{2}+\ConeToneVariable{2}+\ConeToneVariable{3}+\ConeToneVariable{3}+\ConeToneVariable{4}+\ConeToneVariable{4}}{\CsixToneCount{6}}
=\frac{\ConeToneVariable{18}}{\CsixToneCount{6}}
=\CsevenToneSampleMean{3}.
\]

Using $\CsevenToneSampleMean{\bar{x}}=\CsevenToneSampleMean{3}$, the \CsevenToneSampleVariance{sample} \CsevenToneSampleVariance{variance} is

\[
\begin{aligned}
\CsevenToneSampleVariance{s^2}
&=\frac{(\ConeToneVariable{2}\ConeToneVariable{-3})^2+(\ConeToneVariable{2}\ConeToneVariable{-3})^2+(\ConeToneVariable{3}\ConeToneVariable{-3})^2+(\ConeToneVariable{3}\ConeToneVariable{-3})^2+(\ConeToneVariable{4}\ConeToneVariable{-3})^2+(\ConeToneVariable{4}\ConeToneVariable{-3})^2}{\CsixToneCount{6}\CsixToneCount{-1}}\\
&=\frac{\ConeToneVariable{1}+\ConeToneVariable{1}+\ConeToneVariable{0}+\ConeToneVariable{0}+\ConeToneVariable{1}+\ConeToneVariable{1}}{\CsixToneCount{5}}
=\frac45
=\CsevenToneSampleVariance{0.80}.
\end{aligned}
\]

Therefore,

\[
\boxed{\CsevenToneSampleMean{\bar{x}}=\CsevenToneSampleMean{3}=\CfiveTonePopulationMean{\mu},\qquad \CsevenToneSampleVariance{s^2}=\CsevenToneSampleVariance{0.80}\approx\CfiveTonePopulationVariance{\sigma^2}=\CsevenToneSampleVariance{1}.}
\]

\secInicio[sec:c6c7-data-raw][sec:c6-specific-statistics]

\subsection{Daily energy consumption (frequency data)}
\label{sec:c6-probability-frequency}

A facility records daily consumption (in MWh) as follows:

\[
\begin{array}{c|ccc}
	x_i & 8 & 10 & 12 \\ \hline
	f_i & 1 & 2 & 1
\end{array}
\]

The theoretical model has $\mu=10$, $\sigma^2=3$, and
$P(10\leq X\leq12)=0.60$.

\QuestionsTasks[sec:c6-specific-statistics]{
	{6}/{subsec:c6-probability-frequency-b}/
	{Calculate the empirical probability of consumption between $10$ and
	$12$ MWh, inclusive, and its absolute error.
	},
	{7}/{subsec:c7-meanvar-frequency-b}/
	{Calculate the sample mean and sample variance, and compare them with
	the theoretical values.
	}
}

\subsubsection*{C6 Solution}
\phantomsection
\label{subsec:c6-probability-frequency-b}

The total frequency and the frequency in the interval are

\[
\CsixToneCount{n}=\CsixToneCount{1}+2+1=\CsixToneCount{4},
\qquad
\CsixToneCount{k}=\CsixToneCount{2}+1=\CsixToneCount{3}.
\]

Therefore,

\[
\CsixToneEmpirical{\widehat P}(10\leq \ConeToneVariable{X}\leq12)=\frac34=\CsixToneEmpirical{0.75},
\]

and the \CsixToneAbsoluteError{absolute} \CsixToneAbsoluteError{error} is

\[
\CsixToneAbsoluteError{E}=|\CsixToneAbsoluteError{0.75}-0.60|=\CsixToneAbsoluteError{0.15}.
\]

Thus the \CsixToneEmpirical{empirical} \CsixToneEmpirical{probability} is $\ConeToneProbability{0.75}$, which is $0.15$ above the
theoretical probability.

\secInicio[sec:c6-probability-frequency][sec:c6-specific-statistics]

\subsubsection*{C7 Solution}
\phantomsection
\label{subsec:c7-meanvar-frequency-b}

The \CsevenToneSampleMean{sample} \CsevenToneSampleMean{mean} is

\[
\CsevenToneSampleMean{\bar{x}}
=\frac{\ConeToneVariable{1}(\ConeToneVariable{8})+\ConeToneVariable{2}(\ConeToneVariable{10})+\ConeToneVariable{1}(\ConeToneVariable{12})}{\CsixToneCount{4}}
=\frac{\ConeToneVariable{40}}{\CsixToneCount{4}}
=\CsevenToneSampleMean{10}.
\]

The \CsevenToneSampleVariance{sample} \CsevenToneSampleVariance{variance} is

\[
\begin{aligned}
\CsevenToneSampleVariance{s^2}
&=\frac{\ConeToneVariable{1}(\ConeToneVariable{8}\ConeToneVariable{-10})^2+\ConeToneVariable{2}(\ConeToneVariable{10}\ConeToneVariable{-10})^2+\ConeToneVariable{1}(\ConeToneVariable{12}\ConeToneVariable{-10})^2}{\CsixToneCount{4}\CsixToneCount{-1}}\\
&=\frac{\ConeToneVariable{4}+\ConeToneVariable{0}+\ConeToneVariable{4}}{\CsixToneCount{3}}
=\frac83
\approx\CsevenToneSampleVariance{2.67}.
\end{aligned}
\]

Therefore,

\[
\boxed{\CsevenToneSampleMean{\bar{x}}=\CsevenToneSampleMean{10}=\CfiveTonePopulationMean{\mu},\qquad \CsevenToneSampleVariance{s^2}\approx\CsevenToneSampleVariance{2.67}\approx\CfiveTonePopulationVariance{\sigma^2}=\CsevenToneSampleVariance{3}.}
\]

\secInicio[sec:c6-probability-frequency][sec:c6-specific-statistics]

\subsection{Delivery times (grouped data)}
\label{sec:c6-probability-grouped}

A company groups delivery times into three classes:

\[
\begin{array}{c|ccc}
	\text{Delivery time (min)} & 10\text{--}20 & 20\text{--}30 & 30\text{--}40 \\ \hline
	\text{Frequency} & 1 & 2 & 1
\end{array}
\]

The theoretical model has $\mu=25$ minutes, $\sigma^2=75$ minutes$^2$,
and $P(20\leq X<30)=0.45$.

\QuestionsTasks[sec:c6-specific-statistics]{
	{6}/{subsec:c6-probability-grouped-b}/
	{Estimate the empirical probability of a delivery time from $20$ up
	to (but not including) $30$ minutes, and calculate its absolute error.
	},
	{7}/{subsec:c7-meanvar-grouped-b}/
	{Estimate the sample mean and sample variance using class midpoints,
	and compare them with the theoretical values.
	}
}

\subsubsection*{C6 Solution}
\phantomsection
\label{subsec:c6-probability-grouped-b}

The total frequency is $\CsixToneCount{n}=\CsixToneCount{1}+2+1=\CsixToneCount{4}$. The selected class has frequency
$\CsixToneCount{k}=\CsixToneCount{2}$, so

\[
\CsixToneEmpirical{\widehat P}(20\leq \ConeToneVariable{X}<30)=\frac24=\CsixToneEmpirical{0.50}.
\]

The \CsixToneAbsoluteError{absolute} \CsixToneAbsoluteError{error} is

\[
\CsixToneAbsoluteError{E}=|\CsixToneAbsoluteError{0.50}-0.45|=\CsixToneAbsoluteError{0.05}.
\]

Thus the \CsixToneEmpirical{empirical} \CsixToneEmpirical{probability} is $\ConeToneProbability{0.50}$, which is $0.05$ above the
theoretical probability.

\secInicio[sec:c6-probability-grouped][sec:c6-specific-statistics]

\subsubsection*{C7 Solution}
\phantomsection
\label{subsec:c7-meanvar-grouped-b}

The three class midpoints are

\[
15,\qquad25,\qquad35.
\]

Using these midpoints, the estimated \CsevenToneSampleMean{sample} \CsevenToneSampleMean{mean} is

\[
\CsevenToneSampleMean{\bar{x}}
=\frac{\ConeToneVariable{1}(\ConeToneVariable{15})+\ConeToneVariable{2}(\ConeToneVariable{25})+\ConeToneVariable{1}(\ConeToneVariable{35})}{\CsixToneCount{4}}
=\frac{\ConeToneVariable{100}}{\CsixToneCount{4}}
=\CsevenToneSampleMean{25}\text{ min}.
\]

The estimated \CsevenToneSampleVariance{sample} \CsevenToneSampleVariance{variance} is

\[
\begin{aligned}
\CsevenToneSampleVariance{s^2}
&=\frac{\ConeToneVariable{1}(\ConeToneVariable{15}\ConeToneVariable{-25})^2+\ConeToneVariable{2}(\ConeToneVariable{25}\ConeToneVariable{-25})^2+\ConeToneVariable{1}(\ConeToneVariable{35}\ConeToneVariable{-25})^2}{\CsixToneCount{4}\CsixToneCount{-1}}\\
&=\frac{\ConeToneVariable{100}+\ConeToneVariable{0}+\ConeToneVariable{100}}{\CsixToneCount{3}}
=\frac{\ConeToneVariable{200}}{\CsixToneCount{3}}
\approx\CsevenToneSampleVariance{66.67}\text{ min}^2.
\end{aligned}
\]

Therefore,

\[
\boxed{\CsevenToneSampleMean{\bar{x}}=\CsevenToneSampleMean{25}\text{ min}=\CfiveTonePopulationMean{\mu},\qquad
\CsevenToneSampleVariance{s^2}\approx\CsevenToneSampleVariance{66.67}\text{ min}^2\approx\CfiveTonePopulationVariance{\sigma^2}=\CsevenToneSampleVariance{75}\text{ min}^2.}
\]

These grouped-data results are estimates because every observation in a
class is represented by its midpoint.

\secInicio[sec:c6-probability-grouped][sec:c6-specific-statistics]

\subsection{Library visit durations (histogram data)}
\label{sec:c6c7-histogram}

A school librarian records the duration $x$, in minutes, of each visit in
a sample of student visits. The histogram below is the only record provided
of the sample. Each bar shows the number of visits in its class; classes
are left-inclusive and right-exclusive.

\HistogramGraph{(0,10),(10,20),(20,30),(30,40),(40,50)}{2,4,8,4,2}
    {0.75\textwidth}{6cm}{blue!25}{10}{2}

\begin{center}
    Visit duration $x$ (minutes); vertical axis: frequency (number of visits).
\end{center}

\QuestionsTasks[sec:c6-specific-statistics]{
    {6}/{subsec:c6-histogram-solution}/
    {Read each class interval and its frequency from the histogram, and
    determine the total sample size. Calculate the empirical probabilities
    $\widehat P(20\leq X<30)$ and $\widehat P(10\leq X<40)$.
    Find the cumulative frequency and cumulative empirical probability
    below each upper class boundary.
    },
    {7}/{subsec:c7-histogram-solution}/
    {Use class midpoints to estimate the sample mean, sample variance,
    and sample standard deviation. Explain why these statistics are estimates.
    }
}

\subsubsection*{C6 Solution}
\phantomsection
\label{subsec:c6-histogram-solution}

The horizontal axis gives the class boundaries, and the vertical axis has
grid spacing $2$. Reading the bar heights gives

\[
\begin{array}{\ConeToneVariable{c}|\ConeToneVariable{c}\ConeToneVariable{c}\ConeToneVariable{c}\ConeToneVariable{c}\ConeToneVariable{c}}
    \text{Visit duration (min)} & [0,10) & [10,20) & [20,30) & [30,40) & [40,50) \\ \hline
    \CsixToneCount{f_i} & 2 & 4 & 8 & 4 & 2
\end{array}
\]

The total frequency is the \CsixToneCount{sample} \CsixToneCount{size}:

\[
\CsixToneCount{n}=\sum_i \CsixToneCount{f_i}=\CsixToneCount{2}+4+8+4+2=\CsixToneCount{20}.
\]

The class $[20,30)$ has frequency $8$, so

\[
\CsixToneEmpirical{\widehat P}(20\leq \ConeToneVariable{X}<30)=\frac{\ConeToneVariable{8}}{\CsixToneCount{20}}=\CsixToneEmpirical{0.40}.
\]

The interval $[10,40)$ contains three complete classes. Therefore,

\[
\CsixToneCount{k}=\CsixToneCount{4}+8+4=\CsixToneCount{16},
\qquad
\CsixToneEmpirical{\widehat P}(10\leq \ConeToneVariable{X}<40)=\frac{\ConeToneVariable{16}}{\CsixToneCount{20}}=\CsixToneEmpirical{0.80}.
\]

For each upper boundary $\ConeToneVariable{b}$, add the frequencies of all classes below $\ConeToneVariable{b}$.
Divide that cumulative frequency by $\CsixToneCount{n}=\CsixToneCount{20}$ to obtain $\CsixToneEmpirical{\widehat P}(\ConeToneVariable{X}<\ConeToneVariable{b})$:

\[
\begin{array}{\ConeToneVariable{c}|\ConeToneVariable{c}|\ConeToneVariable{c}}
    \ConeToneVariable{b}\text{ (min)} & \text{Cumulative frequency} & \CsixToneEmpirical{\widehat P}(\ConeToneVariable{X}<\ConeToneVariable{b}) \\ \hline
    10 & 2 & 2/20=\ConeToneVariable{0.10} \\
    20 & 2+4=\ConeToneVariable{6} & 6/20=\ConeToneVariable{0.30} \\
    30 & 2+4+8=\ConeToneVariable{14} & 14/20=\ConeToneVariable{0.70} \\
    40 & 2+4+8+4=\ConeToneVariable{18} & 18/20=\ConeToneVariable{0.90} \\
    50 & 2+4+8+4+2=\ConeToneVariable{20} & 20/20=\ConeToneVariable{1}
\end{array}
\]

The last cumulative frequency equals the \CsixToneCount{sample} \CsixToneCount{size}, and the last
cumulative probability is $\ConeToneProbability{1}$, as expected.

\secInicio[sec:c6c7-histogram][sec:c6-specific-statistics]

\subsubsection*{C7 Solution}
\phantomsection
\label{subsec:c7-histogram-solution}

First, calculate each midpoint using $\ConeToneVariable{m_i}=(\ConeToneVariable{a_i}+\ConeToneVariable{b_i})/\ConeToneVariable{2}$. For example,
the midpoint of $[0,10)$ is $(0+10)/2=5$ minutes. The midpoints and
weighted values are

\[
\begin{array}{\ConeToneVariable{c}|\ConeToneVariable{c}\ConeToneVariable{c}\ConeToneVariable{c}\ConeToneVariable{c}\ConeToneVariable{c}}
    \ConeToneVariable{m_i} & 5 & 15 & 25 & 35 & 45 \\ \hline
    \CsixToneCount{f_i} & 2 & 4 & 8 & 4 & 2 \\
    \CsixToneCount{f_i} \ConeToneVariable{m_i} & 10 & 60 & 200 & 140 & 90
\end{array}
\]

The estimated \CsevenToneSampleMean{sample} \CsevenToneSampleMean{mean} is

\[
\CsevenToneSampleMean{\bar{x}}=\frac{\sum_i \CsixToneCount{f_i} \ConeToneVariable{m_i}}{\CsixToneCount{n}}
=\frac{\ConeToneVariable{10}+\ConeToneVariable{60}+\ConeToneVariable{200}+\ConeToneVariable{140}+\ConeToneVariable{90}}{\CsixToneCount{20}}
=\frac{\ConeToneVariable{500}}{\CsixToneCount{20}}=\CsevenToneSampleMean{25}\text{ min}.
\]

Next, use $\CsevenToneSampleMean{\bar{x}}=\CsevenToneSampleMean{25}$ to calculate the weighted squared deviations:

\[
\begin{array}{\ConeToneVariable{c}|\ConeToneVariable{c}\ConeToneVariable{c}\ConeToneVariable{c}\ConeToneVariable{c}\ConeToneVariable{c}}
    \ConeToneVariable{m_i} & 5 & 15 & 25 & 35 & 45 \\ \hline
    \CsixToneCount{f_i} & 2 & 4 & 8 & 4 & 2 \\
    (\ConeToneVariable{m_i}-\CsevenToneSampleMean{\bar{x}})^2 & 400 & 100 & 0 & 100 & 400 \\
    \CsixToneCount{f_i}(\ConeToneVariable{m_i}-\CsevenToneSampleMean{\bar{x}})^2 & 800 & 400 & 0 & 400 & 800
\end{array}
\]

Their sum is $800+400+0+400+800=2400$. Using the sample-variance
denominator $\CsixToneCount{n}-1$, the estimated variance and standard deviation are

\[
\CsevenToneSampleVariance{s^2}=\frac{\sum_i \CsixToneCount{f_i}(\ConeToneVariable{m_i}-\CsevenToneSampleMean{\bar{x}})^2}{\CsixToneCount{n}\CsixToneCount{-1}}
=\frac{\ConeToneVariable{2400}}{\CsixToneCount{20}\CsixToneCount{-1}}=\frac{\ConeToneVariable{2400}}{\CsixToneCount{19}}
\approx\CsevenToneSampleVariance{126.32}\text{ min}^2,
\]
\[
\CsevenToneSampleSD{s}=\sqrt{\CsevenToneSampleVariance{s^2}}=\sqrt{\frac{\ConeToneVariable{2400}}{19}}\approx\CsevenToneSampleSD{11.24}\text{ min}.
\]

Thus the midpoint estimates are

\[
\boxed{\CsevenToneSampleMean{\bar{x}}=\CsevenToneSampleMean{25}\text{ min},\qquad
\CsevenToneSampleVariance{s^2}\approx\CsevenToneSampleVariance{126.32}\text{ min}^2,\qquad \CsevenToneSampleSD{s}\approx\CsevenToneSampleSD{11.24}\text{ min}.}
\]

These statistics are estimates because the histogram does not give the
individual durations: every observation in a class is represented by its
midpoint. The class frequencies and the \CsixToneEmpirical{empirical} \CsixToneEmpirical{probabilities} for unions
of complete classes, however, are read exactly from the histogram.

\secInicio[sec:c6c7-histogram][sec:c6-specific-statistics]

\end{document}
